\section{Related Work}
\label{sec:related-work}

\subsection{Algorithmic stability and generalization}

Classical algorithmic stability studies how a learning rule changes when the
training sample is perturbed, primarily to control generalization error under
an i.i.d.\ sampling model.  The canonical reference is Bousquet and Elisseeff,
who formalize stability notions around replace-one perturbations and fresh-test
evaluation points \cite{bousquet2002stability}.  Kutin and Niyogi expand this
landscape with weaker and almost-everywhere variants
\cite{kutin2002almost}, while Feldman and Vondr\'ak and later Bousquet,
Klochkov, and Zhivotovskiy sharpen high-probability and rate-sensitive uniform
stability bounds \cite{feldman2019high,bousquet2020sharpening}.  Hardt,
Recht, and Singer connect stability to optimization dynamics
\cite{hardt2016train}.  The present paper is orthogonal to this line: our goal
is not a sharper generalization bound but a precise understanding of how
perturbation operations differ at the level of individual finite samples.
Classical stability usually studies replace-one or delete-one perturbations
under expected or fresh-test evaluation, whereas this paper separates
perturbation semantics from evaluation semantics on fixed finite metric
samples.

\subsection{Nearest-neighbor classification and deleted estimates}

The nearest-neighbor line begins with Fix and Hodges and the classical risk
results of Cover and Hart, together with the consistency framework later
formalized by Stone \cite{fix1951discriminatory,cover1967nearest,cover1968estimation,stone1977consistent}.  Those works study statistical performance under
distributional sampling rather than worst-case single-perturbation effects.
Closer to our setting are deleted-estimate and local-rule analyses such as
Rogers and Wagner's finite-sample bound for local discrimination rules and the
distribution-free deleted-estimate inequalities of Devroye and Wagner
\cite{rogers1978finite,devroye1979distribution,devroye1979local}.  These are
important precedents because they show that leave-one-out style evaluation has a
substantial classical history around nearest-neighbor methods.  Even so, their
goal is not to separate LOO from adversarial replace-one at the level of fixed
finite metric witnesses.  Deleted estimates and LOO-style analyses are
historically important for nearest-neighbor rules, but they do not isolate the
fixed-sample adversarial replacement gap studied here.

\subsection{LOO, cross-validation, and conformal stability}

Kearns and Ron connect algorithmic stability with leave-one-out
cross-validation, already emphasizing that a CV-style perturbation has its own
logic \cite{kearns1999algorithmic}.  Recent conformal-prediction work makes
related distinctions operational again, often in service of coverage
certificates rather than 0--1 prediction loss.  Examples include transductive
stability analyses and training-conditional or stable conformal guarantees
\cite{liang2023algorithmic,pournaderi2024training}.  We cite this line because it
reinforces the message that leave-one-out style control and replacement-based
control can lead to different guarantees.  Our setting is more concrete and
more local: deterministic $k$-NN on finite metric spaces with explicit top-$k$
neighbor ordering.  The key message from this literature that carries into our
paper is: LOO is not wrong, it answers a different question.
